% Kun käytät tätä, älä lataa pakettia stfloats tai fix2col. Tämä lataa
% myös paketin fixltx2e.
%% "When the new output routine for LaTeX3 is done, this package will
%% be obsolete. The sooner the better..."
\RequirePackage{dblfloatfix}

% Kun käytät tätä, et enää tarvitse paketteja type1cm ja type1ec:
\RequirePackage{fix-cm}

% Luokkamäärittelyt.
\documentclass[a4paper,onecolumn,11pt,finnish,twoside,final,oldtoc]{boek3}
\usepackage[papersize={185mm,240mm},left=36mm,right=12mm,top=10mm,bottom=25mm,%
            asymmetric,bindingoffset=10mm,marginparwidth=30mm]{geometry}
\reversemarginpar % marginaalit toisin päin
\setcounter{secnumdepth}{7} % kaikki numeroidaan ilman *-optiota

% alleviivaus yms
\usepackage[normalem]{ulem}

% cmap: "Make the PDF files generated by pdflatex "searchable and copyable"
% in Adobe (Acrobat) Reader and other compliant PDF viewers."
\usepackage{cmap}

% "Extended conditional commands"
\usepackage{xifthen}

% Erikoismerkit yms.
\usepackage{textcomp}
\usepackage{texnames}

% AMS & matematiikka
\usepackage{mathtools}
\usepackage[psamsfonts]{amssymb}
\usepackage{amsthm}
\usepackage{latexsym}

% XeLaTeX
\usepackage{fourier}
\usepackage{fontspec}
\defaultfontfeatures{Mapping=tex-text}
\setmainfont[Scale=0.92]{Heuristica}
\usepackage[protrusion=true,verbose=true]{microtype} % XeLaTeX ei toistaiseksi tue 'expansion=true'-optiota

% Otsikoiden säätelyyn
\usepackage{titlesec}

\makeatletter
\g@addto@macro\verbatim{\pdfprotrudechars=0 \pdfadjustspacing=0\relax}
\makeatother

% Babel
\usepackage[finnish,english]{babel}

% if-then-else-rakenteita helposti
\usepackage{ifthen}

% Lisää kokomäärityksiä:
\usepackage[11pt]{moresize}

% Tee jotain jokaisen sivun kohdalla (totpages tarvitsee tätä):
\usepackage{everyshi}

% Laskutehtävien suoritusta (geometry-paketti hyötyy tästä):
\usepackage{calc}

% Riviväli 1,5
\usepackage{setspace}
\onehalfspacing

% URL:ien ladonta ja "tavutus"
\usepackage[obeyspaces,spaces,hyphens,T1]{url}

% TikZ-paketti
\usepackage{pgf,tikz}
\usetikzlibrary{arrows}

% Tunnisteiden ja sivunumeroiden asettamista varten
\usepackage{fancyhdr}

% VerbatimOut Pythonia varten
\usepackage{fancyvrb}
\usepackage{fancybox}

% Euron merkki
\usepackage{eurosym}

% Hymiö
\usepackage{wasysym}

% Harjoitustehtävät answers-paketilla
\usepackage{answers}

% Cancel-paketti supistamista ym. varten
\usepackage{cancel}

% Monipalstainen sisältö esim. lyhyille tehtäville
% Käyttö: \begin{multicols}{n}
%              ...
%         \end{multicols}
\usepackage{multicol}

% Tuki numeroiduille listoille
\usepackage{enumerate}

% Desimaalipilkut
% Käyttö: väli pilkun jälkeen -> listapilkku; ei väliä pilkun jälkeen -> desimaalipilkku
\usepackage{icomma}

% Ei-kursivoidut kreikkalaiset aakkoset, esim. $\upmu$
\usepackage{upgreek}

% Automatisoitu hakemisto
\usepackage{makeidx}
\makeindex

% Aika- ja pvmformaattien kustomointiin
\usepackage{datetime}

% Täältä saa komennon \ifthispageodd
\usepackage{scrextend}

% Marginaalien muuttaminen harjoitustehtäväsivuja varten
\usepackage[strict]{changepage}

% Käännetyt tekstit taulukoita varten
\usepackage{rotating}

\usepackage{array}

% Testailuun, komentoja: \blindtext, \Blindtext, \pagevalues
\usepackage{blindtext}
\usepackage{layouts}

% Ei sisennetä kappaleen ekaa riviä
\setlength{\parindent}{0.0cm}

\theoremstyle{definition}
\newtheorem{theorem}{Teoreema}

% Lisäattribuutteja \includegraphics-komentoon
\usepackage{graphicx}

% allaolevia loogisia muuttujia voi muuttaa meta.tex-tiedostossa
% esim. \varitfalse

% väreissä vai mustavalkoisena
% vaatii vastaavat kuvatiedostot
\newif\ifvarit
\varittrue

% versionumero kansilehteen
\newif\ifversio
\versiotrue

% mikrorahoituskanavat kirjan alkuun
\newif\ifmikro
\mikrotrue

\newcommand{\kurssinTunnus}{}
\newcommand{\kurssinNimi}{}
\newcommand{\sitaatti}{Sitaatti puuttuu.}
\newcommand{\sitaatinLahde}{}
\newcommand{\kirjanVersio}{}
\newcommand{\kirjanPainopaikka}{}
\newcommand{\kirjanVastaavuus}{}
\newcommand{\metasivu}{}

% Hypersetup (metatason pdf-säätöä)
\input{content/hypersetup}

% Omat komennot
% pitkän matematiikan kursseja
\newcommand{\maaI}{Funktiot ja yhtälöt}
\newcommand{\maaII}{Polynomifunktiot}
\newcommand{\maaIII}{Geometria}
\newcommand{\maaIV}{Analyyttinen geometria}
\newcommand{\maaV}{Vektorit}
\newcommand{\maaVI}{Todennäköisyys ja tilastot}
\newcommand{\maaVII}{Derivaatta}
\newcommand{\maaVIII}{Juuri- ja logaritmifunktiot}
\newcommand{\maaIX}{Trigonometriset funktiot ja lukujonot}
\newcommand{\maaX}{Integraalilaskenta}
\newcommand{\maaXI}{Lukuteoria ja logiikka}
\newcommand{\maaXII}{Numeerisia ja algebrallisia menetelmiä}
\newcommand{\maaXIII}{Differentiaali- ja integraalilaskennan jatkokurssi}
\newcommand{\maaXIV}{Kertaus}

% lyhyen matematiikan kursseja
\newcommand{\mabI}{Lausekkeet ja yhtälöt}
\newcommand{\mabII}{Geometria}
\newcommand{\mabIII}{Matemaattisia malleja I}
\newcommand{\mabIV}{Matemaattinen analyysi}
\newcommand{\mabV}{Tilastot ja todennäköisyys}
\newcommand{\mabVI}{Matemaattisia malleja II}
\newcommand{\mabVII}{Talousmatematiikka}
\newcommand{\mabVIII}{Matemaattisia malleja III}
\newcommand{\mabIX}{Kertaus}

% nuolia
\newcommand{\ekvi}{\Longleftrightarrow}
\newcommand{\impli}{\Longrightarrow}

% sanoja math-moodissa
\newcommand{\kun}{\textnormal{kun} \;}
\newcommand{\jos}{\textnormal{jos} \;}
\newcommand{\tai}{\textnormal{tai} \;}

% yleisimpiä joukkoja
\newcommand{\N}{\mathbb{N}}
\newcommand{\Z}{\mathbb{Z}}
\newcommand{\Q}{\mathbb{Q}}
\newcommand{\R}{\mathbb{R}}
\newcommand{\C}{\mathbb{C}}

\input{content/meta}

\ifvarit
    \newcommand{\vari}{}
    \definecolor{vapaa_matikka_vari_3}{cmyk}{0.2,0.27,0.0,0.0}	% Laatikot
    \definecolor{vapaa_matikka_vari_4}{cmyk}{1.0,0.42,0.65,0.0}	% Esimerkit
    \definecolor{vapaa_matikka_vari_5}{cmyk}{0.60,0.80,0.0,0.0}	% Tehtävänumerot
\else
    \newcommand{\vari}{_grey}
    \definecolor{vapaa_matikka_vari_3}{cmyk}{0.0,0.0,0.0,0.25}	% Laatikot
    \definecolor{vapaa_matikka_vari_4}{cmyk}{0.0,0.0,0.0,0.6}	% Esimerkit
    \definecolor{vapaa_matikka_vari_5}{cmyk}{0.0,0.0,0.0,0.8}	% Tehtävänumerot
\fi

% Tarvitaan kuvien ja taulukkojen vierekkäin laittamiseen.
\def\vcent#1{\mathsurround0pt$\vcenter{\hbox{#1}}$}

\usepackage{../commons/packages/esimerkki}
\usepackage{../commons/packages/kansilehti}
\usepackage{../commons/packages/kuva}
\usepackage{../commons/packages/kuvaaja}
\usepackage{../commons/packages/laatikot}
\usepackage{../commons/packages/logiikka}
\usepackage{../commons/packages/lukusuora}
\usepackage{../commons/packages/merkkikaavio}
\usepackage{../commons/packages/otsikkotyylit}
\usepackage{../commons/packages/ppalkki}
\usepackage{../commons/packages/qrlinkki}
\usepackage{../commons/packages/tehtavat}
\usepackage{../commons/packages/termi}
\usepackage{../commons/packages/todistus}
\usepackage{../commons/packages/tunnisteet}
\usepackage{../commons/packages/valit}


\let\cleardoublepage\clearpage

% Lisäfontteja
% Thin - ExtraLight - Light - Regular - Medium - SemiBold - Bold - ExtraBold - Black
\newfontface\fontI[Scale=3.1]{Raleway-Light.otf}
\newfontface\fontII[Scale=2.3]{Raleway-ExtraLight.otf}
\newfontface\fontIII[Scale=1.5]{Raleway-Light.otf}
\newfontface\fontIIIc[Scale=1.5,Color=vapaa_matikka_vari_5]{Raleway-Light.otf}
\newfontface\fontIV[Scale=0.9]{Raleway-Medium.otf}
\newfontface\fontIVc[Scale=0.9,Color=vapaa_matikka_vari_5]{Raleway-Medium.otf}
\newfontface\fontV[Scale=1.2]{Raleway-Regular.otf}

\begin{document}

\urlstyle{same}

\newdateformat{finshortdate}{\THEDAY.\THEMONTH.\THEYEAR}
\selectlanguage{finnish}
\finshortdate

\ifversio
    \kansilehti{Versio \kirjanVersio \, (\today)}
\else
    \kansilehti{}
\fi

\newpage
\subsection*{Lisenssi \includegraphics[scale=0.3]{../commons/logo/cc.pdf} \includegraphics[scale=0.3]{../commons/logo/by.pdf}}

Tämän teoksen käyttöoikeutta koskee Creative Commons Nimeä 3.0 Muokkaamaton (CC BY 3.0) -lisenssi.
This work is licensed under a Creative Commons Attribution 3.0 Unported (CC BY 3.0) License.

\qrlinkki{http://creativecommons.org/licenses/by/3.0/deed.fi}{Lisenssin suomenkielinen tiivistelmä}
\qrlinkki{http://creativecommons.org/licenses/by/3.0/deed.en}{Lisenssin englanninkielinen tiivistelmä}
\qrlinkki{http://creativecommons.org/licenses/by/3.0/legalcode}{Tarkat lisenssiehdot (englanniksi)}

Sinulla on vapaus:
\begin{itemize}
\item jakaa — kopioida, levittää, näyttää ja esittää teosta
\item remiksata — valmistaa muutettuja teoksia
\item käyttää teosta kaupallisiin tarkoituksiin
\end{itemize}
Seuraavilla ehdoilla:
\begin{itemize}
\item nimeä — Teoksen tekijä on ilmoitettava siten kuin tekijä tai teoksen lisensoija on sen määrännyt (mutta ei siten että ilmoitus viittaisi lisenssinantajan tukevan lisenssinsaajaa tai Teoksen käyttötapaa)
\end{itemize}

Määräys lisäyksenä lisenssiin: kaikkien kirjoittajien nimet on ilmoitettava jossakin kohtaa kirjaa.

\newpage
\input{content/people}

\ifmikro
    \newpage \subsection*{Mikrolahjoituskanavat}

\qrlinkki{https://flattr.com/profile/oppikirjamaraton}{Flattr}

\qrlinkki{bitcoin:148pMeTViRMFBqMVZRnMZ6Hwccq9WTubq1?label=Oppikirjamaraton}{Bitcoin}

\fi

\newpage
\subsection*{Uusimmat versiot}

\textbf{Git-repot} \\
\url{https://github.com/avoimet-oppimateriaalit-ry}

\textbf{Valmiit pdf-versiot} \\
\url{http://goo.gl/7OikD}

\subsection*{Projekti}

\textbf{Tavoite} \\
Kokonainen lukion pitkän matematiikan kirjasarja CC BY -lisenssillä

\textbf{Koordinaattori} \\
Vesa Linja-aho (Metropolia Ammattikorkeakoulu)

\textbf{Rahoitus} \\
Teknologiateollisuuden 100-vuotissäätiö: 30~000 \euro \; (2013)

\subsection*{Avoimet oppimateriaalit ry}

\textbf{Kotisivu} \\
\url{http://avoimetoppimateriaalit.fi}

\textbf{Wiki (beta)} \\
\url{http://koti.kapsi.fi/~kenkku/moin.fcgi}

\textbf{Sähköpostit} \\
\href{mailto:yhdistys@avoimetoppimateriaalit.fi}{yhdistys@avoimetoppimateriaalit.fi} \\
\href{mailto:vapaamatikka@avoimetoppimateriaalit.fi}{vapaamatikka@avoimetoppimateriaalit.fi}

\textbf{Facebookissa} \\
\url{https://facebook.com/avoimetoppimateriaalit} \\
\url{https://facebook.com/oppikirjamaraton}

\textbf{IRCnetissä} \\
\#avoimetoppimateriaalit

\tableofcontents
\addtocontents{toc}{\vspace*{-50pt}} % move start of toc higher on page
\input{content/content} % incl. Opensolutionfile, Closesolutionfile, appendices

\newpage
\addcontentsline{toc}{chapter}{Hakemisto}
\printindex

\end{document}
